%% Supplemental material for "Physics-Exact Credit Assignment for Multi-Agent
%% Flutter Suppression" (submitted to the Journal of Guidance, Control, and
%% Dynamics). Cited from the main text; not subject to the journal's page
%% limit. Reproduces exactly the content of the corresponding main-text
%% sections in the prior (longer) draft, unabridged, for reviewers and readers
%% who want the full derivations and reproducibility detail.
\documentclass[11pt]{article}
\usepackage[T1]{fontenc}
\usepackage{lmodern}
\usepackage[letterpaper,margin=1in]{geometry}
\usepackage{amsmath,amssymb,amsthm}
\usepackage{array}
\usepackage{booktabs}
\usepackage{graphicx}
\usepackage{siunitx}
\usepackage[hidelinks]{hyperref}
\usepackage[round]{natbib}

\newtheorem{proposition}{Proposition}
\newtheorem{remark}{Remark}

\newcommand{\Uf}{U_{\mathrm{f}}}
% Generated by scripts/fill_numbers.py -- do not edit numeric macros by hand.
\newcommand{\FlutterSpeed}{137.35}
\newcommand{\FeasibleCells}{6}
\newcommand{\TotalCells}{9}
\newcommand{\LocalFbStabilised}{4}
\newcommand{\LocalFbHealthy}{-0.03}
\newcommand{\LqrFixedStabilised}{5}
\newcommand{\LqrFixedHealthy}{-5.15}
\newcommand{\LqrSchedStabilised}{5}
\newcommand{\LqrSchedHealthy}{-5.08}
\newcommand{\LqgLocalStabilised}{5}
\newcommand{\LqgLocalHealthy}{-5.52}
\newcommand{\LqrOracleStabilised}{6}
\newcommand{\LqrOracleHealthy}{-5.07}
\newcommand{\CommOracleStabilised}{6}
\newcommand{\CommOracleHealthy}{-0.18}
\newcommand{\IppoPhysicsStabilised}{6}
\newcommand{\IppoPhysicsHealthy}{-0.25}
\newcommand{\MoccaStabilised}{5}
\newcommand{\MoccaHealthy}{-0.14}
\newcommand{\RecurrentOnlyStabilised}{5}
\newcommand{\RecurrentOnlyHealthy}{-0.04}
\newcommand{\BestRobustVariant}{comm\_oracle}
\newcommand{\BestRobustStabilised}{6}
\newcommand{\BestNominalVariant}{ippo\_physics}
\newcommand{\BestNominalHealthy}{-0.25}
\newcommand{\ResultsSummarySentence}{\emph{[ResultsSummarySentence: to be written]}}
\newcommand{\ResultsNarrative}{\emph{[ResultsNarrative: to be written]}}
\newcommand{\ResultsAblations}{\emph{[ResultsAblations: to be written]}}
\newcommand{\PhaseNarrative}{\emph{[PhaseNarrative: to be written]}}
\newcommand{\ConclusionSentence}{\emph{[ConclusionSentence: to be written]}}


\title{Supplemental Material: Physics-Exact Credit Assignment for\\
Multi-Agent Flutter Suppression}
\author{}
\date{}

\begin{document}
\maketitle

\noindent This document collects material referenced from the main text but
moved here to meet the journal's page-length guidance. Section, table, and
figure numbers below are internal to this document (prefixed S) and are
cross-referenced from the main text by name, not by number. All numerical
results here were computed by the same code and are covered by the same
seed-count and statistical-correction standard stated where noted in the main
text.

\section*{S1. Training hyperparameters}
\label{app:hyperparameters}

Table~\ref{tab:params-ppo} lists the PPO and network hyperparameters shared
across every learned variant in the credit and architecture ablations and
the exploration-scale sweep, except where a hyperparameter is itself the
subject of an ablation (the initial exploration standard deviation, and the
presence or absence of the architectural mechanisms). All
runs share parameters across agents, with the agent index supplied as a one-hot
input, so the parameter count does not grow with the number of surfaces.

\begin{table}[h]
\centering
\caption{PPO and network hyperparameters.}
\label{tab:params-ppo}
\begin{tabular}{@{}lp{5.6cm}@{}}
\toprule
Hyperparameter & Value \\
\midrule
Parallel environments & 128 \\
Rollout length & 128 steps \\
Epochs per update & 4 \\
Minibatches per epoch & 4 \\
BPTT segment length & 32 steps \\
Learning rate (Adam) & $3\times 10^{-4}$ \\
Discount $\gamma$ & 0.995 \\
GAE $\lambda$ & 0.95 \\
PPO clip range & 0.2 \\
Value loss coefficient & 0.5 \\
Entropy coefficient & $2\times 10^{-3}$ (0 in the exploration sweep) \\
Max gradient norm & 0.5 \\
Curriculum fraction & 0.4 of training \\
Hidden layer width (trunk) & 128 \\
Phasor encoder hidden width & 64 \\
Retained modes per phasor & 2 \\
Consensus rounds & 2 \\
Action-head output-layer gain (init.) & 0.01 \\
\bottomrule
\end{tabular}
\end{table}

\section*{S2. Compute and reproducibility}
\label{app:compute}

Every learned variant in this paper trains for $1.5{\times}10^6$ environment
steps across 128 vectorised environments (Table~\ref{tab:params-ppo}) on a
single NVIDIA RTX 5000 Ada (laptop) GPU, using PyTorch 2.6. Across the 86
training runs behind the main results and ablations, median wall-clock training
time is \SI{331}{\second} (minimum \SI{246}{\second} for a plain policy,
maximum \SI{1115}{\second} for the full \texttt{mocca} architecture, whose
GRU encoder and consensus rounds are the most expensive forward pass in the
ablation set), at a median throughput of
$4.5{\times}10^3$ environment steps per second --- the batched environment
rather than the network is the bottleneck at this
scale, since the plant assembly is vectorised across
environments but not across agents. The full set of runs reported in this paper --- the credit and architecture
ablations, the eight-surface scaling probe, and every additional probe in
Section~S3 below --- totals \SI{12.7}{\hour} of GPU time; no result in
this paper required infrastructure beyond a single workstation GPU.

\section*{S3. Additional probes}
\label{sec:probes}

The results in the main text are held to a fixed standard: six seeds,
Welch's $t$-test, Cohen's $d$ and Hedges' $g$, Holm-Bonferroni corrected.
Three further questions arose directly out of those results and were
checked, but only at one or two seeds each --- not because the questions
matter less, but because reaching the main standard for every question this
study raised would have multiplied its cost several times over. We report
them as probes rather than findings: Table~\ref{tab:probes} gives every seed
individually rather than a mean, so that no reader mistakes a two-point
spread for an estimated standard deviation.

\begin{table}[h]
\centering
\caption{Additional probes, one or two seeds each, at $\sigma=0.05$. Every
healthy-decay value is listed per seed, not averaged, because two points do not
estimate a standard deviation. \texttt{-{}-} marks a probe not run for that
variant.}
\label{tab:probes}
\small
\begin{tabular}{@{}>{\raggedright\arraybackslash}p{3.6cm}l>{\raggedright\arraybackslash}p{3.4cm}c@{}}
\toprule
Probe & Variant & Healthy decay [1/s], per seed & Held \\
\midrule
\multicolumn{4}{@{}p{0.97\linewidth}@{}}{\emph{Do the negative mechanisms of the architecture ablation also hurt as a residual correction?}} \\
residual + raw comm. & \texttt{residual\_comm} & $-0.697,\ -0.247$ & 6/9, 6/9 \\
residual + health-channel consensus & \texttt{residual\_health} & $-0.518,\ -0.066$ & 5/9, 5/9 \\
\addlinespace
\multicolumn{4}{@{}p{0.97\linewidth}@{}}{\emph{Does supervising the encoder on the true modal phasor help?}} \\
supervised encoder, plain head & \texttt{aux\_encoder} & $-0.219,\ -0.029$ & 5/9, 5/9 \\
supervised encoder, residual & \texttt{aux\_residual} & $-0.431,\ -0.510$ & 5/9, 5/9 \\
\addlinespace
\multicolumn{4}{@{}p{0.97\linewidth}@{}}{\emph{Eight-surface scaling}} \\
consensus, plain head, $N{=}8$ & \texttt{mocca\_nohead} & $-0.028,\ -0.015$ & 9/9, 8/9 \\
recurrent only, $N{=}8$ & \texttt{recurrent\_only} & $+0.135,\ -0.011$ & 9/9, 9/9 \\
\bottomrule
\end{tabular}
\end{table}

\paragraph{The negative architecture result survives moving to a residual policy}
The plain-policy architecture ablation found that consensus
and raw neighbour sharing hurt. Applying the same two mechanisms to a policy
that corrects a decentralised LQG base law instead of replacing it --- the
setting where \texttt{residual\_only} is the strongest result in the paper at
$-1.53\,\mathrm{s}^{-1}$ --- gives
\texttt{residual\_comm} at $-0.697$ and $-0.247$ and \texttt{residual\_health}
at $-0.518$ and $-0.066$: both clearly below \texttt{residual\_only}, at both
seeds. The negative result is not an artefact of the specific plain-policy
setting it was first measured in.

\paragraph{Supervising the encoder on ground truth does not help}
\texttt{aux\_encoder} adds a loss term supervising the phasor encoder directly
on the true modal phasor --- privileged information, unavailable outside
simulation --- rather than letting the encoder's representation be shaped
purely by the RL objective. Both seeds ($-0.219$, $-0.029$) sit below the
un-supervised \texttt{mocca\_nohead} result of $-0.308\pm0.169$, and the
residual counterpart
\texttt{aux\_residual} ($-0.431$, $-0.510$) sits far below
\texttt{residual\_only}'s $-1.53$. We read this as mild evidence that the
credit signal already supplies what the encoder needs to represent, and that
an additional supervised loss competes with, rather than complements, the
policy gradient here --- consistent with the broader pattern that mechanisms
motivated by the physics did not survive contact with this specific plant
and reward.

\paragraph{Eight surfaces changes what ``held'' means, not just how many}
The main text reports that all nine evaluation conditions become
feasible for the learned policies at eight surfaces; Table~\ref{tab:probes}
gives the numbers behind that sentence. They deserve a more careful reading
than ``held'' alone provides: \texttt{mocca\_nohead} holds nine of nine and
eight of nine conditions, but at a healthy-condition decay rate of only
$-0.028$ and $-0.015$ --- an order of magnitude weaker than the three-surface
result for the same architecture. \texttt{recurrent\_only} holds every
condition at both seeds, yet one seed's healthy-condition decay is
\emph{positive} ($+0.135$). ``Held'' in Table~\ref{tab:probes} means the
episode did not cross the divergence threshold within its duration, not that
the wing was driven to a decisively damped state; a bounded, barely-growing
oscillation can hold a condition by this criterion without being a controller
we would recommend. We stand by the qualitative claim that eight-surface
architectures fail less catastrophically than three-surface ones, and retract
any stronger reading of it than that.

\section*{S4. Baseline controller synthesis}
\label{app:baselines}

This section gives the synthesis equations behind every rung of the
classical baseline ladder in the main text, so that ``discrete-time LQG on
local accelerometers'' is reproducible from this document alone rather than
only from the repository.

\subsection*{Augmented, discretised plant}

Actuator lag is part of the design plant, not an unmodelled disturbance added
afterwards: with $z=[q,\dot q, x_{\mathrm{w}}, \delta]$ (the state of the
equation of motion augmented with the deflections) and first-order servo lag of
time constant $\tau_k$ per surface,
\begin{equation}
\dot z = \tilde A z + \tilde B u, \qquad
\tilde A = \begin{bmatrix} A & B \\ 0 & -\mathrm{diag}(\tau_k^{-1}) \end{bmatrix},
\quad
\tilde B = \begin{bmatrix} 0 \\ \mathrm{diag}(\tau_k^{-1}) \end{bmatrix},
\label{eq:augplant}
\end{equation}
where $A,B$ are the plant matrices in state-space form and
$u\in[-1,1]^K$ is the normalised command. Every controller in this study is
synthesised in discrete time at the \SI{200}{\hertz} control rate via the exact
zero-order-hold discretisation
\begin{equation}
\begin{bmatrix} \Phi & \Gamma \\ 0 & I \end{bmatrix}
= \exp\!\left( \begin{bmatrix} \tilde A & \tilde B \\ 0 & 0 \end{bmatrix} \Delta t \right),
\label{eq:zoh}
\end{equation}
giving $z_{t+1} = \Phi z_t + \Gamma u_t$. Designing in continuous time and
integrating the resulting gain with a forward-Euler step at $\Delta t=\SI{5}{ms}$
is not a minor approximation here: the wake states have $|\lambda|$ up to
$\sim\!\SI{1.3e4}{\radian\per\second}$, so $\Delta t\,|\lambda|\approx 63$, and
the observer built that way is unstable regardless of how good its gain is.

\subsection*{LQR gain}

The state cost weights physical structural energy rather than an arbitrary
diagonal, plus a light penalty on held deflection so the optimum does not park
surfaces at the travel limit:
\begin{equation}
Q = \begin{bmatrix} K_{\mathrm{modal}} & & \\ & M_{\mathrm{modal}} & \\ & & \mathrm{diag}\!\left(w_{\mathrm{rate}}/\delta_{\max,k}^2\right) \end{bmatrix},
\qquad
R = \mathrm{diag}\!\left(w_{\mathrm{effort}}/\delta_{\max,k}^2\right),
\label{eq:lqrcost}
\end{equation}
with $w_{\mathrm{effort}}=2{\times}10^4$, $w_{\mathrm{rate}}=10^2$ throughout.
The discrete-time algebraic Riccati equation
\begin{equation}
P = \Phi^{\!\top} P \Phi - \Phi^{\!\top} P \Gamma \left(R+\Gamma^{\!\top} P \Gamma\right)^{-1} \Gamma^{\!\top} P \Phi + Q
\label{eq:dare}
\end{equation}
gives the gain $\mathcal{K}=(R+\Gamma^{\!\top} P \Gamma)^{-1}\Gamma^{\!\top} P \Phi$
and command $u_t=-\mathcal{K} z_t$ \citep{anderson2007optimal}.
\texttt{BatchedLQR} solves \eqref{eq:dare} once at a fixed design speed;
\texttt{BatchedScheduledLQR} solves it on a \SI{9}{}-point airspeed grid over
the evaluation envelope and linearly interpolates $\mathcal{K}$ between grid
points; \texttt{BatchedJamAwareLQR} additionally zeros the failed surface's
column of $\Gamma$ per failure hypothesis and resolves \eqref{eq:dare} for each
hypothesis on the same grid, so the surviving surfaces' gains already account
for the loss rather than compensating for it reactively.

\subsection*{Local rate feedback}

The floor of the ladder uses no plant model: with $\dot h_k,\dot\alpha_k$ the
locally sensed plunge and twist rate at surface $k$'s own station,
\begin{equation}
u_k = -\left(g_h\,\dot h_k + g_\alpha\,\dot\alpha_k\right), \qquad g_h=0.35,\ g_\alpha=0.9,
\label{eq:ratefeedback}
\end{equation}
clipped to $[-1,1]$.

\subsection*{Steady-state Kalman observer on local accelerometers}

The centralised and decentralised LQG rungs replace the true state $z_t$ in
$u_t=-\mathcal{K}\hat z_t$ with an estimate driven only by the accelerometer
channels the agents themselves observe, $y_t = H z_t$, where $H$ is built from
the acceleration rows of the plant rather than a selection matrix, since an
accelerometer measures $\ddot q$, an algebraic function of the state rather
than a state itself. The steady-state predictor-form filter is
\begin{align}
\Sigma &= \Phi\Sigma\Phi^{\!\top} + W
- \Phi\Sigma H^{\!\top}\!\left(H\Sigma H^{\!\top}+V\right)^{-1}\! H\Sigma\Phi^{\!\top},
\label{eq:kalmanriccati} \\
L &= \Phi\Sigma H^{\!\top}\!\left(H\Sigma H^{\!\top}+V\right)^{-1}, \\
\hat z_{t+1} &= \Phi\hat z_t + \Gamma u_t + L\left(y_t - H\hat z_t\right),
\label{eq:kalmanupdate}
\end{align}
with process noise $W=10^{-4}I$ and measurement noise
$V=10^{-6}\,\mathrm{diag}(\lVert H_i\rVert^2)$ scaled per-channel, since
accelerometer rows of $H$ have norms of order $10^4$--$10^5$ and an unscaled
$V$ demands impossible measurement precision and drives $L$ to blow up
\citep{anderson2007optimal}. \texttt{BatchedLQG} solves
\eqref{eq:kalmanriccati}--\eqref{eq:kalmanupdate} once with $H$ built from all
six channels and $\mathcal{K}$ commanding all three surfaces jointly, on the
same speed grid as the scheduled LQR. \texttt{Batched\-Decentralized\-LQG}
solves an \emph{independent} instance of
\eqref{eq:kalmanriccati}--\eqref{eq:kalmanupdate} per surface $k$, with $H$
restricted to $k$'s own two accelerometer rows and $\Gamma$ restricted to $k$'s
own column, so agent $k$'s estimate, gain and command never depend on any
other surface's measurement or actuation --- the discrete-time, model-based
analogue of the information structure the learned policies are given. Its
scheduling grid spans the same $[1.0,1.5]\,\Uf$ range as the training
curriculum exactly, by construction rather than
coincidence, so the comparison in the main text is not confounded
by the classical design being scheduled over a narrower envelope than the
learned policy was trained on.

\subsection*{Re-tuning the decentralised LQG for control-authority margin}
\label{app:robust}

\texttt{BatchedRobustDecentralizedLQG} is
structurally identical to \texttt{BatchedDecentralizedLQG}: same per-surface
observer, same restriction to each agent's own two accelerometers and own
actuator. Two remedies for its lack of a guaranteed margin
\citep{doyle1978guaranteed} were checked. The first, loop transfer recovery
\citep{doylestein1981ltr}, boosts the process-noise weighting $W$ in
\eqref{eq:kalmanriccati} by $\rho\, b_k b_k^{\top}$ (agent $k$'s own reduced
input column) as $\rho\to\infty$, pushing the filter loop toward the
full-state loop shape. Sweeping $\rho$ from $0$ to $10^4$ changes the
control-authority margin (the largest reduction in true control authority,
relative to the design assumption, the closed loop remains stable under, found
by bisection on the discrete-time closed-loop spectral radius) not at all: it
stays at $0.158$ at $1.45\,\Uf$ throughout. Comparing against the same gain
applied with perfect (observer-free) state feedback, whose margin is also
$0.158$, shows why: the observer is not the bottleneck, so there is nothing
for an observer-side remedy to recover.

The second remedy, and the one \texttt{BatchedRobustDecentralizedLQG} actually
uses, re-weights the LQR cost \eqref{eq:lqrcost} itself: $w_{\mathrm{effort}}$
increased from $2\times10^4$ to $10^6$, with $w_{\mathrm{rate}}$, $W$ and $V$
unchanged. This raises the control-authority margin from $0.158$ to $0.179$
before plateauing (no further gain to $w_{\mathrm{effort}}=10^{7}$), at a cost
to nominal decay rate documented in the main text.
\texttt{scripts/robust\_baseline\_check.py} reproduces both checks.

\bibliographystyle{plainnat}
\bibliography{refs,refs_extra}
\end{document}

